\chapter{Text Formatting}

OmniLaTeX ships sensible defaults for text formatting so that most
documents look polished out of the box.  This chapter covers the
micro-typographic engine, dash handling, quotation marks, and spacing
controls.

% =========================================================================
\section{Microtype}
% =========================================================================

The \texttt{microtype} package improves typographic quality through three
mechanisms that are all enabled by default in OmniLaTeX:

\begin{description}
  \item[Protrusion] Characters such as commas, periods, hyphens, and
    quotation marks are allowed to extend slightly into the right margin,
    eliminating visible raggedness caused by punctuation at line ends.
  \item[Expansion] Inter-character spacing is microscopically adjusted
    (typically $\pm 1$\,--\,$2$\,emu) to achieve even line lengths
    without hyphenation.
  \item[Tracking] Global letter-spacing is modified slightly at line
    beginnings and endings to balance paragraph shape.
\end{description}

No configuration is required.  To disable microtype globally (not
recommended), pass \texttt{microtype=false} to the document class.

% =========================================================================
\section{Extended Dashes}
% =========================================================================

The \texttt{extdash} package provides shorthand macros that produce
correctly kerned en-dashes and em-dashes without relying on \TeX's
ligature-based \texttt{-{}-} / \texttt{-{}-{}-} mechanism.  This is
especially useful in documents that use hyphens inside compound words
alongside ranges.

\begin{table}[htbp]
  \centering
  \caption{extdash shortcuts.}
  \label{tab:extdash}
  \begin{tabular}{@{}lll@{}}
    \toprule
    Input                       & Output            & Meaning \\
    \midrule
    \texttt{fast\textbackslash-/paced} & fast\-/paced & Em-dash joining hyphenated words \\
    \texttt{1\textbackslash--10}       & 1\--10       & En-dash range \\
    \texttt{\textbackslash--{}-}       & ---          & Em-dash \\
    \texttt{\textbackslash-}           & -            & Discretionary hyphen (no break) \\
    \bottomrule
  \end{tabular}
\end{table}

Key advantage: \texttt{fast\-/paced} produces \emph{fast}—\emph{paced}
(an em-dash) whereas the ligature approach \texttt{fast---paced} would
incorrectly typeset as fast---paced.

% =========================================================================
\section{Quotations}
% =========================================================================

The \texttt{csquotes} package provides language-aware quotation marks via
\cs{enquote}\texttt{\{text\}}.  OmniLaTeX loads \texttt{csquotes} by
default and integrates it with the bibliography:

\begin{verbatim}
\enquote{Hello world}
% → "Hello world"  (English)
% → ,Hello world'  (British)
% → »Hello world«  (German)
\end{verbatim}

Nested quotes are handled automatically:

\begin{verbatim}
\enquote{She said, \enquote{Hello world}.}
% → "She said, 'Hello world'."
\end{verbatim}

To use \cs{enquote} with \texttt{biblatex}, set:

\begin{verbatim}
\SetCiteCommand{\parencite}
\end{verbatim}

This ensures that \cs{textcite} and \cs{parencite} produce properly
quoted author names in the running text.

% =========================================================================
\section{Ragged Typesetting}
% =========================================================================

The \texttt{ragged2e} package improves upon plain \cs{raggedright} by
still hyphenating words when necessary, producing a cleaner right margin
than \LaTeX's built-in ragged mode.  OmniLaTeX loads \texttt{ragged2e}
and exposes two commands:

\begin{table}[htbp]
  \centering
  \caption{Ragged typesetting commands.}
  \label{tab:ragged}
  \begin{tabular}{@{}lp{0.6\textwidth}@{}}
    \toprule
    Command          & Behaviour \\
    \midrule
    \texttt{\textbackslash RaggedRight} & Left-aligns with hyphenation; ideal for
      narrow columns, posters, and slide content. \\
    \texttt{\textbackslash RaggedLeft}  & Right-aligns with hyphenation; useful for
      captions or marginal notes. \\
    \bottomrule
  \end{tabular}
\end{table}

Use \cs{RaggedRight} instead of \cs{raggedright} inside
\texttt{tabular}, \texttt{minipage}, and \texttt{tcolorbox} environments
to avoid overly loose inter-word spacing.

% =========================================================================
\section{Blind Text}
% =========================================================================

The \texttt{blindtext} package generates filler paragraphs and complete
dummy documents for layout testing:

\begin{verbatim}
\Blindtext[3][1]      % 3 paragraphs, 1-level sectioning
\Blinddocument        % Full document with chapters and sections
\end{verbatim}

OmniLaTeX sets the blind-text language to match \cs{mainlang}, so
\cs{Blindtext} produces English filler text by default and will switch
to German, French, etc.\ when the document language is changed.

% =========================================================================
\section{Line Spacing}
% =========================================================================

OmniLaTeX uses \texttt{setspaceenhanced} for fine-grained line-spacing
control.  The default baseline skip factor is:

\begin{verbatim}
\linespread{1.05}
\end{verbatim}

This produces slightly tighter lines than \LaTeX's default $1.2$ factor
while remaining readable for body text.  For wider spacing:

\begin{verbatim}
\onehalfspacing     % 1.5× line spacing
\doublespacing      % 2.0× line spacing
\end{verbatim}

All spacing commands affect the current scope, so they can be used inside
groups to localise changes:

\begin{verbatim}
{\onehalfspacing
  This paragraph has wider spacing.
}
\end{verbatim}

% =========================================================================
\section{Paragraph Spacing}
% =========================================================================

KOMA-Script provides the \texttt{parskip} option to control vertical
spacing between paragraphs instead of \TeX's default indentation.
OmniLaTeX configures this via \cs{documentparspacing}:

\begin{verbatim}
\documentparspacing{half}   % half-line skip, no indent (default for manual)
\documentparspacing{full}   % full-line skip, no indent
\documentparspacing{indent} % traditional indent, no skip
\end{verbatim}

The \texttt{half} setting is the default for the \texttt{manual} document
type and produces comfortable visual separation between paragraphs without
wasting vertical space.  The \texttt{indent} setting is used for the
\texttt{thesis} and \texttt{article} types.
